\documentclass{itsas_pfc}
\usepackage[utf8]{inputenc}
\usepackage[T1]{fontenc}
\usepackage{geometry}
\usepackage{graphicx}
\usepackage{float}
\usepackage{hyperref}
\usepackage[english,indonesian]{babel}

\title{Optimasi Pemrosesan Template LaTeX Menggunakan Pendekatan Berbasis AI}
\author{Academic Support Agent}
\date{\today}

\begin{document}
\maketitle

\begin{abstract}

\end{abstract}

\chapter{PENDAHULUAN}


\chapter{TINJAUAN PUSTAKA}


% Original content (commented out):
% %%%%%%%%%%%%%%%%%%%%%%%%%%%%%%%%%%%%%%%%%%%%%%%%%%%%%%%%
%                                                      %
%  Plantilla para Perfil de Proyecto de Investigación  %
%                                                      %
%%%%%%%%%%%%%%%%%%%%%%%%%%%%%%%%%%%%%%%%%%%%%%%%%%%%%%%%
%
% CONFIG: Definición del estilo del documento
%
\documentclass[a4paper,12pt, oneside]{itsas_pfc}
%
% CONFIG: Fichero con configuración
%
\input{Config/pfc_options.tex}

%estableciendo el estilo de la bibliografia APA
\usepackage[backend=biber,style=apa]{biblatex}
%Agregando la fuente bibliografica
\addbibresource{bibliografia/ejemplosApaVersion07.bib}
% Comienza el documento en sí
\begin{document}
%
% CONFIG: dar valores a ciertas variables
%
\input{Config/datosTesistas.tex}
\input{Config/datosTesis.tex}
%         
% CONTENIDO: Página de portada
%         
\input{Config/cover.tex}
%
% CONFIG: Definir el estilo de las cabeceras/piés de página.
%
%\input{Config/headers.tex}
%
% CONFIG: Estilo de los capítulos
%
\input{Config/chapters.tex}
%
% CONTENIDO: Tras esto puedes incluir todos los capítulos/secciones que desees.
%
\input{contenido/ArchivoContenido} % Capítulo 1
%
%
% CONFIG: Estilo de la bibliografía
%
\input{Config/bibliography.tex}
%
%
%\clearpage
%\input{anexos/ArchivoAnexo} % Apéndice I
%
% Aquí concluye todo
%

\end{document}


\end{document}